%! Parametrically Defined Circles and Lines — Unit 4, Lesson 4.4 — Lesson Plan
\documentclass[10pt]{article}
\usepackage{precalculus-article}
\usepackage{precalculus-boxes}
\usepackage{graphicx}
\graphicspath{{images/}}
\newcommand{\TallMath}[1]{$\displaystyle #1\rule[-1.4em]{0pt}{3.2em}$}
\newcommand{\CourseName}{AP Precalculus}
\newcommand{\SchoolYear}{2026--2027}
\newcommand{\MeetingLength}{55 minutes}
\newcommand{\UnitNumberName}{Unit 4: Functions Involving Parameters, Vectors, and Matrices \quad}
\newcommand{\LessonNumberName}{Lesson 4.4: Parametrically Defined Circles and Lines}

\begin{document}

%----------------------------------------------------------------------
% Title Block
%----------------------------------------------------------------------
\begin{center}
  {\Large\bfseries \CourseName: \SchoolYear} \\[4pt]
  {\small\bfseries \UnitNumberName \LessonNumberName \quad (2 instructional periods, \MeetingLength\ each)}
\end{center}

%----------------------------------------------------------------------
% Primary Objective
%----------------------------------------------------------------------
\begin{tcolorbox}[colback=sky, colframe=navy, boxrule=0.9pt, arc=2mm,
    left=2mm, right=2mm, top=1.5mm, bottom=1.5mm]
  \textbf{Primary Objective:} Students will express circular paths
  parametrically using transformations of $(x(t),\,y(t)) = (\cos t,\,\sin t)$
  and model linear motion along a line segment using an initial point and
  constant rates of change. \textit{(Big Idea: Functions)}
\end{tcolorbox}

\vspace{0.08in}

%----------------------------------------------------------------------
% Priority Ideas & Skills
%----------------------------------------------------------------------
\begin{skillbox}[Priority Ideas \& Skills]{goldbox}
\begin{tabularx}{\linewidth}{@{} p{0.46\linewidth} X @{}}
\textbf{AP Skills} & \textbf{Key Understandings (EKs)} \\
\midrule
\textbf{1.B} \textit{Express in equivalent forms} ---
  Rewrite a verbal or geometric description of circular or linear
  motion as a pair of parametric equations, and identify the radius,
  center, direction, and speed from the algebraic form.
&
\textbf{EK 4.4.A.1}: A complete counterclockwise revolution around
the unit circle that starts and ends at $(1,\,0)$, centered at the
origin, is modeled by $(x(t),\,y(t)) = (\cos t,\,\sin t)$ with
domain $0 \le t \le 2\pi$. \\[4pt]

\textbf{1.C} \textit{Construct new functions via transformations} ---
  Apply horizontal/vertical shifts and scale factors to $(\cos t,\,\sin t)$
  to obtain the parametrization of any circle; modify the argument of
  $t$ to control direction and speed.
&
\textbf{EK 4.4.A.2}: Transformations of $(\cos t,\,\sin t)$ can model
any circular path: scaling $r$ sets the radius, shifting by $(h,\,k)$
moves the center, and replacing $t$ with $-t$ reverses orientation. \\[4pt]

&
\textbf{EK 4.4.A.3}: A linear path along the segment from
$(x_1,\,y_1)$ to $(x_2,\,y_2)$ can be parametrized by
$x(t)=x_1+(x_2-x_1)t$, $y(t)=y_1+(y_2-y_1)t$ for $0\le t\le 1$,
using an initial position and constant rates of change. \\
\end{tabularx}
\end{skillbox}

\vspace{0.08in}

%----------------------------------------------------------------------
% Vocabulary
%----------------------------------------------------------------------
\begin{skillbox}[{Vocabulary, Concepts, \& Theorems}]{greenbox}
\begin{tabularx}{\linewidth}{@{} p{0.26\linewidth} X @{}}
\textbf{Term} & \textbf{Definition / Note} \\
\midrule
unit circle parametrization &
  $(x(t),\,y(t))=(\cos t,\,\sin t)$, $0\le t\le 2\pi$; models one
  full counterclockwise revolution around the unit circle. \\[4pt]
radius / center &
  In $(h+r\cos t,\;k+r\sin t)$, $r$ is the radius and $(h,\,k)$
  is the center of the circular path. \\[4pt]
orientation &
  The direction of travel along a curve: \emph{counterclockwise}
  when $t$ increases in $({\cos t},\,{\sin t})$; \emph{clockwise}
  when $t$ is replaced by $-t$. \\[4pt]
linear parametrization &
  $x(t)=x_1+(x_2-x_1)t$, $y(t)=y_1+(y_2-y_1)t$, $0\le t\le 1$;
  traces the segment from $(x_1,\,y_1)$ to $(x_2,\,y_2)$ at
  constant speed. The coefficients $(x_2-x_1)$ and $(y_2-y_1)$
  are the rates of change of $x$ and $y$ with respect to $t$. \\
\end{tabularx}
\end{skillbox}

\vspace{0.08in}

%----------------------------------------------------------------------
% Activate Prior Knowledge
%----------------------------------------------------------------------
\begin{skillbox}[Activate Prior Knowledge \& Spiral Review]{sky}
\begin{tabularx}{\linewidth}{@{}X c@{}}
\begin{minipage}[t]{0.96\linewidth}\vspace{0pt}
\textbf{Spiral topics activated during Warm-Up (first 8 minutes):}
\begin{itemize}[leftmargin=1.4em, itemsep=2pt]
  \item \textbf{Cosine and sine values}: $\cos$ and $\sin$ at $0$, $\pi/2$,
    $\pi$, $3\pi/2$, $2\pi$ --- the key outputs that define the unit circle path.
  \item \textbf{Function transformations}: vertical/horizontal shifts and
    scale factors applied to sine and cosine --- maps directly to the
    $(h + r\cos t,\; k + r\sin t)$ template.
  \item \textbf{Counterclockwise motion on the unit circle}: interpreting
    quarter-turn and half-turn positions geometrically --- connects the angle
    $t$ to a physical location on the circle.
\end{itemize}
\end{minipage}
&
\end{tabularx}
\end{skillbox}

\vspace{0.08in}

%----------------------------------------------------------------------
% Hook
%----------------------------------------------------------------------
\begin{skillbox}[Hook --- Entry Event]{sky}
\textbf{Scenario:} A Ferris wheel has a radius of 30 ft and its center
is 35 ft above the ground. A rider boards at the rightmost seat, which
starts at ground height $+$ 35 ft, exactly 30 ft to the right of center.

\textbf{Discussion prompts:}
\begin{enumerate}[leftmargin=1.4em, itemsep=3pt]
  \item \textit{``Write parametric equations that give the rider's height
    $y$ and horizontal position $x$ at time $t$ as the wheel completes
    one counterclockwise revolution.''}
  \item \textit{``What if the wheel runs clockwise instead? How do you
    change the equations?''}
  \item \textit{``What is the rider's highest point, and what value of $t$
    achieves it?''}
\end{enumerate}

\textbf{Key tension:} A description like ``height varies as a sine wave''
captures $y$ alone. The parametric form tracks \emph{both} coordinates
simultaneously and encodes direction --- essential for real motion.
\end{skillbox}

\vspace{0.08in}

%----------------------------------------------------------------------
% Lesson
%----------------------------------------------------------------------
\begin{skillbox}[Lesson --- Instruction Sequence]{sky}
\begin{multicols}{2}

\textbf{Step 1 --- Unit Circle Parametrization} (8 min)

Write $(x(t),\,y(t)) = (\cos t,\,\sin t)$, domain $0 \le t \le 2\pi$.
Evaluate at $t = 0, \pi/2, \pi, 3\pi/2, 2\pi$ to fill a class table.
Identify the starting point $(1,\,0)$, the quarter-turn point $(0,\,1)$,
the half-turn point $(-1,\,0)$, and the return to $(1,\,0)$.
Emphasize: \emph{counterclockwise} direction, one full revolution.
(EK 4.4.A.1)

\columnbreak

\textbf{Step 2 --- Transformations of the Circle} (12 min)

General form: $x(t) = h + r\cos t$, $y(t) = k + r\sin t$.
Show how each parameter changes the picture:
\begin{itemize}[leftmargin=1.2em, itemsep=2pt]
  \item $r$ scales the radius.
  \item $(h,k)$ shifts the center.
  \item Replacing $t$ with $-t$: $\cos(-t)=\cos t$, $\sin(-t)=-\sin t$
    --- reverses orientation to \emph{clockwise}.
  \item Replacing $t$ with $ct$ ($c > 1$) speeds up the revolution;
    replacing with $t/(2\pi) \cdot T$ scales the period to $T$.
\end{itemize}
Apply to Ferris wheel: $h=0$, $k=35$, $r=30$.
(EK 4.4.A.2)

\end{multicols}

\begin{multicols}{2}

\textbf{Step 3 --- Linear Parametrization} (10 min)

Segment from $(x_1,\,y_1)$ to $(x_2,\,y_2)$:
\[
  x(t) = x_1 + (x_2-x_1)\,t,\quad y(t) = y_1 + (y_2-y_1)\,t,\quad 0\le t\le 1.
\]
At $t=0$: start point. At $t=1$: end point. At $t=0.5$: midpoint.
The coefficients are rates of change: $\Delta x / \Delta t$ and $\Delta y / \Delta t$.
Connect to slope: $\Delta y / \Delta x = (y_2-y_1)/(x_2-x_1)$.
(EK 4.4.A.3)

\columnbreak

\textbf{Step 4 --- Apply Both Contexts} (10 min)

\textbf{Circular:} Ferris wheel --- state equations, find highest point
($y$ is max when $\sin t = 1$, i.e.\ $t = \pi/2$; height $= 65$ ft).
For one revolution in 40 seconds, scale: $t_{\text{real}} = 40t/(2\pi)$,
so rider is highest at $t_{\text{real}} = 10$ seconds.

\textbf{Linear:} Delivery drone flies from $(2,\,3)$ to $(8,\,11)$.
Equations: $x(t) = 2 + 6t$, $y(t) = 3 + 8t$, $0\le t\le 1$.
Find position at $t=0.5$: $(5, 7)$ --- the midpoint.

\end{multicols}

\smallskip
\textbf{Pacing note:} Steps 1--2 fill Period 1; Steps 3--4 + activity in Period 2.
\end{skillbox}

\vspace{0.08in}

%----------------------------------------------------------------------
% Active Monitoring
%----------------------------------------------------------------------
\begin{skillbox}[Active Monitoring]{redbox}
\begin{itemize}[leftmargin=1.4em, itemsep=3pt]
  \item \textbf{Radius vs.\ center confusion}: Students may swap $r$
    and $(h,k)$ in the formula. Prompt: \textit{``Which part of
    $h + r\cos t$ shifts the whole circle, and which part controls how
    wide it is?''}
  \item \textbf{Clockwise sign error}: Students may negate both
    components. Prompt: \textit{``What is $\cos(-t)$? What is
    $\sin(-t)$? So which one changes sign when we reverse direction?''}
  \item \textbf{Linear endpoints}: Students may evaluate $t=0$ and $t=1$
    and get the wrong direction. Ask: \textit{``Where does your curve
    start? Where does it end? Does that match the problem?''}
  \item \textbf{Highest point on a circle}: Students may guess the center
    height. Prompt: \textit{``What value of $\sin t$ is the largest
    possible? What $t$ gives that value?''}
\end{itemize}
\end{skillbox}

\vspace{0.08in}

%----------------------------------------------------------------------
% Group Work & Differentiation
%----------------------------------------------------------------------
\begin{skillbox}[Group Work \& Differentiation]{redbox}
Students work in groups of 3--4 on the tiered activity.

\begin{itemize}[leftmargin=1.4em, itemsep=4pt]
  \item \textbf{Tier R (Remediate)}: Write the parametrization for a
    circle centered at the origin with radius 5 starting at $(5,\,0)$.
    Verify that the point at $t = \pi/2$ is $(0,\,5)$.
  \item \textbf{Tier A (Approaching Proficiency)}: Write the
    parametrization for a circle centered at $(-2,\,4)$ with radius 3
    traversed \emph{clockwise}. Also write the parametrization for
    the segment from $(1,\,2)$ to $(7,\,-4)$.
  \item \textbf{Tier E (Extension)}: For the Ferris wheel, determine
    the parameter value $t$ when the rider is at the highest point.
    If the wheel completes one revolution in 40 seconds, at what
    actual time (in seconds) is the rider highest? Generalize: if the
    period is $T$ seconds, write a formula for $t_{\text{real}}$ as a
    function of $T$.
\end{itemize}
\end{skillbox}

\vspace{0.08in}

%----------------------------------------------------------------------
% Individual Work & Assessment
%----------------------------------------------------------------------
\begin{skillbox}[Individual Work \& Assessment]{redbox}
Exit ticket administered independently (final 5 minutes of Period 2).

\begin{enumerate}[leftmargin=1.4em, itemsep=4pt]
  \item A circle is centered at $(3,\,-1)$ with radius $4$. Write
    parametric equations for counterclockwise motion starting at $(7,\,-1)$.
  \item Write a parametrization for the line segment from $(0,\,6)$ to
    $(4,\,0)$. At what parameter value $t$ does the segment cross the
    $x$-axis?
\end{enumerate}

\textbf{Data use:} Students who cannot write the Ferris-wheel equations
without scaffolding need more support with the transformation template;
revisit in the opening of Lesson 4.5.
\end{skillbox}

\vspace{0.08in}

%----------------------------------------------------------------------
% Reinforcement & Extension
%----------------------------------------------------------------------
\begin{skillbox}[Reinforcement \& Extension]{goldbox}
\textbf{Homework (Lesson 4.4):} 5--6 problems covering: writing
circular parametrizations from a description, evaluating at specific
$t$-values, finding when a circular path reaches a given point,
writing and evaluating a linear segment parametrization, and one
AP-style multiple-choice item.

\textbf{Preview of Lesson 4.5 --- Implicitly Defined Functions:}
Ask students: \textit{``We wrote $x(t) = h + r\cos t$ and
$y(t) = k + r\sin t$. If you eliminate $t$ using the identity
$\cos^2 t + \sin^2 t = 1$, what equation in $x$ and $y$ do you get?''}
Lesson 4.5 explores how to recover the Cartesian equation from a
parametric form and what is gained or lost.
\end{skillbox}

\end{document}
